\section{Threats to Validity\label{sec:threats}}

\parg{Generalizability:} We evaluate ten open-source systems in Java, C, and C++, each with one reference architecture. The findings may not extend to other systems or alternative architectural decompositions. The study is transductive: representation learning and graph processing can use unlabeled test nodes and their edges. It therefore does not establish performance on newly arriving files outside the observed graph or on maintaining a mapping as the system evolves. Other information sources, including change history and dynamic dependencies, are outside the evaluated scope.

\parg{Dependencies and reference mappings:} Dependencies are extracted using static analysis, which may miss or misresolve dependencies. Preprocessing and build configuration can also affect the dependencies available for analysis. Because predictors that use neighbors depend on the extracted graph, such errors can affect their predictions. Models with neighbors enabled and disabled are trained separately. Their comparison therefore measures the effect of training with access to neighbor information, rather than removing neighbors from a fixed trained model.

We made a best-effort attempt to identify and obtain the source versions corresponding to the published mappings, using the available version information and documentation. Despite these efforts, some mismatches between the analyzed file sets and the mappings remain. These may reflect incomplete documentation of the original source versions or build configurations, which limits our ability to establish an exact correspondence between the analyzed source code and the published mappings.

Files can plausibly belong to more than one module, and annotators may assign them differently. Because the initial labeled set is sampled from the same mapping used for evaluation, an incorrect assignment can both mislead the model during training and penalize an architecturally plausible prediction.

\parg{Experimental variation:} We use 100 runs per configuration. These runs reuse the same systems and can include the same files, so they are not independent samples. With this many runs, small differences can reach statistical significance, so we report results per system and include effect sizes to distinguish the magnitude of a difference from its significance. The comparisons were not fixed in advance, so the analyses are exploratory and may favor patterns specific to this benchmark.

The runs use different initial labeled sets, but the contribution of stochastic model training to the observed variation was not assessed separately. Input representations were generated once per system and reused across runs, so variation arising from their generation was not assessed. The seeds governing model initialization and training were not retained, limiting exact reproducibility. Cliff's $\delta$ treats the paired scores as two independent groups, so it describes the separation of the two score distributions and not the paired differences.

\parg{Interpretation of results:} Performance measures agreement with the reference architecture, not whether a prediction is the only architecturally reasonable assignment. Folder organization may reflect the same decisions used to define that architecture, so strong directory-based predictions do not by themselves demonstrate recovery of architectural intent. The directory-error analysis applies only where the directory baseline makes errors, which excludes C.Img, SH-3D, and T.Mates in most or all splits, and it does not isolate files whose module differs from the majority module of their directory.
Message-passing depth was varied only for HGT in RQ2, with one and two layers, so the results do not show how deeper models or other predictors respond. It reports accuracy only, so it does not show how performance differs between small and large modules.

Relationships between performance and structural properties are reported as associations, not causal effects, since the benchmark does not allow these properties to be separated from one another or from system size and label availability. The fine-tuning gain is that of the whole dependency-guided procedure, since no alternative adaptation objective was evaluated, so it is not attributed to dependency topology alone.
